% !TEX root = ../main.tex
% =====================================================================
% 7. EXPERIMENTAL SETUP AND ENGINEERING     target: 2.5 - 3 pages
%
% "How I did it". Two things live here:
%   (a) the scientific protocol (data, baselines, metrics, seeds),
%   (b) the engineering I built to make the study possible.
% The examiner said NOT to give a full dataset description or a full
% hyperparameter list -- keep those short, and push details to the
% appendix.
% =====================================================================

\section{Experimental Setup}
\label{sec:setup}

\begin{guide}
Keep this section factual and compact. A reader should be able to answer: what data,
what baselines, what recipe, what metrics, how many runs. Anything longer than three
lines about hyperparameters belongs in \Cref{app:repro}.
\end{guide}

% ---------------------------------------------------------------------
\subsection{Datasets}
\label{sec:setup:data}

\begin{guide}
Half a page, no more. One paragraph per dataset family, and one small table. Do not
list the 128 \ucr{} datasets here -- the appendix can hold the per-dataset table.
\textbf{Questions to answer:} Why is \webtraffic{} special (it is the only one with
per-time-step ground truth, so it is the only place NDCG@$n$ can be computed)? Why
also run 85 \ucr{} datasets (to show the result is not one lucky dataset)?
\end{guide}

\paragraph{\webtraffic{}.}
\todo{Three or four sentences: synthetic dataset from \citet{early2024millet},
\num{n} series of length \num{T}, \num{C} classes, and -- the important part -- the
time steps that were modified are known, so explanations can be scored directly. Say
that I used it as my screening set: fast to train, and the only place explanation
quality can be measured against the truth.}

{\webtraffic{} is a synthetic dataset introduced by \citet{early2024millet} for exactly
this purpose: unlike a real dataset, the generator knows which time steps it modified
to create each class, so those steps can be used as ground truth to score an
explanation instead of only trusting it by eye. It has 1{,}000 series in total (500 for
training, 500 for test), each of length 1{,}008, split across 10 classes. Because it is
small and fast to train (a few minutes per model on a GPU), I used it as my screening
set: every new encoder or pooling head is tried here first, and only the ones that look
promising are then run on the much slower \ucr{} sweep.}

\paragraph{\ucr{} archive.}
\todo{Two or three sentences: the \ucr{} archive \citep{dau2019ucr} holds 128
univariate datasets of very different lengths and sizes; \millet{} published results
on 85 of them, so I sweep the same 85 for a fair head-to-head, plus the rest for
completeness.}

{The \ucr{} archive \citep{dau2019ucr} holds 128 univariate time series datasets, from
tiny (16 training series) to large (almost 9{,}000), and short (15 steps) to long
(nearly 2{,}900 steps). \citet{early2024millet} published \millet{} results on 85 of
these 128, so I run every model on that same 85 first, which is what makes a
dataset-by-dataset comparison against their numbers possible. I also run the remaining
43 datasets, so the headline result is not resting on one lucky subset.}

\begin{table}[h]
\caption{The two dataset families used in this study. \todo{Fill in the numbers from
\texttt{results/SEA\_NET/data\_summary.csv}.}}
\label{tab:datasets}
\begin{center}
\small
\begin{tabular}{lrrrl}
\toprule
\textbf{Dataset} & \textbf{\#series} & \textbf{Length} & \textbf{\#classes} &
\textbf{Per-step ground truth?} \\
\midrule
\webtraffic{}       & 1{,}000 (500 / 500) & 1{,}008 & 10 & yes \\
\ucr{} (85 published) & 40--16{,}637 & 24--2{,}709 & 2--60 & no \\
\ucr{} (128 total)    & 40--24{,}000 & 15--2{,}844 & 2--60 & no \\
\bottomrule
\end{tabular}
\end{center}
\end{table}
% Numbers checked against results/SEA_NET/data_summary.csv (train+test count per
% dataset) and results/SEA_NET/seanet_conjunctive.../comparison_vs_millet.csv (the 85
% rows that have a non-empty millet_acc are the 85 MILLET published; the other 43 are
% the rest of the archive).

% ---------------------------------------------------------------------
\subsection{Baselines}
\label{sec:setup:baselines}

\begin{guide}
Be very clear about one point, because it is the strongest methodological decision of
the internship: I did not only compare against numbers printed in a paper, I
\emph{re-trained} the published backbones myself with my own recipe. Explain why that
matters (different recipes give different numbers, and the comparison would otherwise
be unfair in my favour or against me).
\end{guide}

\todo{One paragraph: FCN \citep{wang2017fcn}, ResNet \citep{wang2017fcn} and
InceptionTime \citep{fawaz2020inceptiontime}, each with \millet{} Conjunctive pooling,
trained here with exactly the same recipe as \seanet{}. Also report the published
\millet{} numbers, but never mix them with my own in the same column.}

{I compare against three backbones: FCN, ResNet \citep{wang2017fcn}, and InceptionTime
\citep{fawaz2020inceptiontime}, the same one \millet{} itself uses. All three are paired
with \millet{}'s own Conjunctive pooling head (\texttt{mil\_conjunctive}), and all three
are trained through the exact same pipeline and recipe as \seanet{} (\Cref{sec:setup:protocol}) --
same optimiser, same epoch budget, same early stopping, same data splits. This matters
because a paper's published number is a mix of the architecture \emph{and} whatever
recipe that paper's authors used; if I only quoted the published number I could not
tell whether a difference in my results comes from the architecture or from a different
number of training epochs. Re-training the baselines myself removes that confound: any
gap between \seanet{} and a re-trained baseline is attributable to the architecture,
because everything else is held fixed. I confirmed the recipe genuinely matters and is
not free to ignore: using \millet{}'s own published hyperparameters instead of mine
(\texttt{configs/models/sv1/millet\_paper.yaml} -- 1500 epochs, no weight decay, no
label smoothing, no validation split) gives a different accuracy again. For this
reason the published \millet{} numbers are always kept in their own column and never
averaged or merged with a number I produced myself.}

% ---------------------------------------------------------------------
\subsection{Training protocol}
\label{sec:setup:protocol}

\begin{guide}
Short paragraph plus a compact table. The examiner does not want every hyperparameter,
so give the ones that affect a comparison (optimiser, epochs, batch size, early
stopping, seeds) and send the rest to the appendix.
\textbf{Questions to answer:} What is fixed for every model, and why does that make
the comparison fair? How many seeds, and why that number? How is early stopping done
when a dataset has very few training series?
\end{guide}

\todo{One paragraph describing the recipe in words, then:}

{Every model in this study -- \seanet{} and every baseline -- is trained with the same
recipe (\texttt{seanet/train.py}), so the numbers in \Cref{sec:results} differ only
because of the model, not the training setup. I optimise a label-smoothed
cross-entropy loss with Adam, for up to 400 epochs, and stop early once 60 epochs pass
with no improvement in the loss I am monitoring. When a dataset has at least 100
training series I hold out a stratified 20\% validation split and monitor validation
loss; on the smaller \ucr{} datasets a 20\% split would be a handful of series and too
noisy to trust, so I monitor training loss instead. The batch size scales with the
dataset (a tenth of the training set, capped at 16, floor of 2), because a fixed batch
size of 16 is more than half the whole training set on the smallest \ucr{} datasets.}

\begin{table}[h]
\caption{Training recipe, identical for every model in this study.
\todo{Fill from \texttt{configs/main.yaml}.}}
\label{tab:protocol}
\begin{center}
\small
\begin{tabular}{ll}
\toprule
\textbf{Setting} & \textbf{Value} \\
\midrule
Optimiser              & Adam \citep{kingma2015adam}, lr $=1.25\times10^{-3}$, weight decay $=1.2\times10^{-4}$ \\
Epochs / early stopping& max 400 epochs; stop after 60 epochs with no improvement \\
Batch size             & $\min(\lfloor n_{\text{train}}/10 \rfloor, 16)$, floor of 2 \\
Validation split       & stratified 80/20 when $n_{\text{train}} \geq 100$; else train-loss early stopping \\
Label smoothing        & 0.13 \\
Seeds                  & 1 (seed $=0$); no repeats -- see note below \\
Hardware               & Grid'5000 GPU nodes (Lille, Sophia) \\
\bottomrule
\end{tabular}
\end{center}
\end{table}

\note{I checked this against the actual results files: every row of
\texttt{results/SEA\_NET/<model>/results.csv} has \texttt{seed=0}, and there is no
second seed anywhere in the 66-model sweep. The table's original draft said ``3 seeds,
mean $\pm$ std'' -- that was the plan, but with 66 model variants $\times$ 129 datasets
already close to the internship's compute budget, running each 3$\times$ was not
possible in the time available. This is a real limitation: single-seed accuracy
numbers, especially on the smallest \ucr{} datasets (some have only 16 training
series), can move by a few points on a different seed. I say this plainly in
\Cref{sec:discussion} rather than implying a variance estimate I do not have.}

% ---------------------------------------------------------------------
\subsection{Metrics}
\label{sec:setup:metrics}

\todo{Short paragraph: accuracy and cross-entropy loss for classification; AOPCR and
NDCG@$n$ for explanation quality (defined in \Cref{sec:bg:metrics}); parameters,
model size, FLOPs, training time and inference latency for cost. Say explicitly that
reporting cost prevents a model from ``winning'' just by being bigger.}

{I report three groups of numbers for every model. \emph{Classification}: test
accuracy and cross-entropy loss. \emph{Explanation quality}: AOPCR and NDCG@$n$
(\Cref{sec:bg:metrics}), computed only on \webtraffic{} for NDCG since it is the only
dataset with per-step ground truth. \emph{Cost}: parameter count and saved model size
for every model, plus FLOPs, training time and inference latency for the ones profiled
in \texttt{results/SEA\_NET/profile.csv}. Reporting cost alongside accuracy matters
because a bigger model can win on accuracy by brute force alone; without the parameter
column, a reader cannot tell whether a gain came from a better idea or just from more
capacity.}

\begin{guide}
Add one warning sentence here about AOPCR being unnormalised, so a reader does not
compare my AOPCR column with the published one. Repeating this once in
\Cref{sec:results} is fine -- it is the single easiest number to misread.
\end{guide}

\note{AOPCR is unnormalised (it depends on the model's own logit scale), so my AOPCR
numbers are only meaningful \emph{relative to each other}, never as an absolute value
against a number from another paper -- including \millet{}'s own published AOPCR. The
\texttt{millet\_aopcr} column in my results tables is kept only as a sanity check that
my implementation of the metric matches theirs on their own model, not as a number to
beat directly.}

% ---------------------------------------------------------------------
\subsection{The experimental pipeline I built}
\label{sec:setup:pipeline}

\begin{contribution}
The whole experimental pipeline -- configuration system, model registry, resumable
sweeps, results store, figure and table generation, experiment tracking -- is my own
work. It is what made it possible to train and compare 66 model variants over
129 datasets within the internship.
\end{contribution}

\begin{guide}
This subsection is where the software engineering side of the internship becomes
visible, and examiners value it. Do not turn it into a code tour: give the design
idea, one diagram, and the benefit of each choice in half a sentence.
\textbf{Questions to answer:} Why a config-driven design instead of editing code for
each experiment? How does a sweep survive a job that is killed by the scheduler? How
do I guarantee that no number in this report was typed by hand?
\end{guide}

\todo{One paragraph: a model is always \emph{encoder + pooling head}, both chosen by
name in a YAML file, so any combination can be trained without touching the code. Each
model writes into its own results folder, one row per dataset, updated in place.}

{A model is always \emph{encoder + pooling head}, and both are chosen by name in a
config file (\texttt{configs/models/<model>.yaml}) rather than hardcoded. Two small
registries (\texttt{seanet/features.py} for encoders, \texttt{seanet/pooling.py} for
pooling heads) map a name like \texttt{sea\_mstcn\_sep\_gated} or
\texttt{sea\_topk\_conjunctive} to the matching class; \texttt{build\_model\_from\_config}
(\texttt{seanet/model.py}) looks up both names and wires them together
(\Cref{sec:architecture}). This is what let 66 model variants exist without 66 copies
of near-identical training code: every new row in the results tables of
\Cref{sec:results} is a new YAML file, not a new script. Each model also writes into its own results folder under
\texttt{results/SEA\_NET/<model>/}, one row per dataset in \texttt{results.csv}, so two
models' numbers can never overwrite each other by accident.}

\todo{One paragraph: sweeps are resumable (a finished dataset is recorded and skipped
on restart), which matters because Grid'5000 jobs have a time limit. Runs are tracked
with MLflow, and every figure and table in this report is regenerated by a single
command from the stored results.}

{Grid'5000 jobs have a wall-clock limit, and a full sweep over 129 datasets for one
model does not always finish inside it. Every model tracks which datasets it has
already finished in a small \texttt{done\_train\_dataset.txt} file; restarting the same
sweep command skips anything already done and picks up where it stopped, so a killed
job costs at most the one dataset that was mid-training, not the whole sweep. Every
run is also logged to MLflow (\texttt{mlflow.db}), recording its hyperparameters,
metrics and the trained weights, so any run can be inspected or reloaded later without
re-training it. Finally, every figure and \LaTeX{} table used in this report is
generated from the stored \texttt{results.csv} files by a single command
(\texttt{python main.py paper}), which is what lets me say no number in this report
was typed in by hand.}

\begin{figure}[t]
  \centering
  \figph{\textbf{Figure X -- Training and evaluation pipeline.} Config (YAML)
  $\rightarrow$ data loading $\rightarrow$ model build (encoder + pooling head)
  $\rightarrow$ training with early stopping $\rightarrow$ evaluation (accuracy, loss,
  AOPCR, NDCG) $\rightarrow$ per-model results store $\rightarrow$ automatic figures
  and \LaTeX{} tables $\rightarrow$ this report. Show the resume loop and the MLflow
  side branch.}
  % ---- Diagram brief for fig:pipeline (draw this yourself) ----
  % One flowchart, left to right, with two things branching off the main line:
  %
  % [1] configs/main.yaml (which model + which dataset)
  %      + configs/models/<model>.yaml (encoder name + pooling name + training block)
  %       |
  %       v
  % [2] seanet/data.py: D.load_dataset(name) -> train/test series
  %      -> apply_instance_representation() (single_point, unchanged; or sliding_window, cut into windows)
  %       |
  %       v
  % [3] seanet/model.py: build_model_from_config()
  %      -> build_encoder(cfg.encoder)   [seanet/features.py registry]
  %      -> build_pooling(cfg.pooling)   [seanet/pooling.py registry]
  %      -> EncoderPoolNet(encoder, pool)
  %       |
  %       v
  % [4] seanet/train.py: SeaNetModel.fit(train_ds, val_ds)
  %      loop over epochs: forward -> label-smoothed CE + lambda_entropy * attention_entropy(attn)
  %      -> backward -> Adam step; monitor val (or train) loss; stop after `patience` epochs flat
  %       BRANCH UP into a small box: "resume check" -- before [4] starts, main.py checks
  %       done_train_dataset.txt for this model; if this dataset is already listed, skip
  %       straight to the next dataset instead of re-training (this is the loop that survives
  %       a job killed by the Grid'5000 scheduler).
  %       |
  %       v
  % [5] seanet/train.py: score_model()
  %      -> safe_evaluate() -> test accuracy, balanced accuracy, AUROC, loss
  %      -> model.evaluate_interpretability() -> AOPCR (always), NDCG@n (WebTraffic only)
  %       |
  %       v
  % [6] one row appended to results/SEA_NET/<model>/results.csv (dataset, seed, params,
  %      accuracy, loss, AOPCR, NDCG, train_time_s, ...)
  %       BRANCH DOWN into a side box: "MLflow" -- the same row's metrics, plus the run's
  %      hyperparameters and (optionally) the trained weights, are logged to mlflow.db via
  %      seanet/tracking.py, in parallel with [6], not instead of it.
  %       |
  %       v
  % [7] python main.py leaderboard / paper / results
  %      -> reads every model's results.csv -> results/SEA_NET/leaderboard.csv
  %      -> results/paper_figures/*.pdf (figures) and results/paper_figures/tables/*.tex (tables)
  %       |
  %       v
  % [8] \input{} / \includegraphics{} -> this report (main.tex)
  %
  % Draw [1]-[8] as one left-to-right chain; [4]'s resume-check box sits just above it with an
  % arrow looping back to "next dataset"; [6]'s MLflow box sits just below it as a side branch,
  % not a step the main chain waits on.
  \caption{\todo{Caption: the path from a config file to a number in this report; no
  value is copied by hand.}}
  \label{fig:pipeline}
\end{figure}

{The path from \texttt{main.py} to a number in this report has no manual step: every
figure and every table cell is produced by the pipeline in \Cref{fig:pipeline}, so a
number cannot be copy-pasted wrong or left stale after a re-run.}

\todo{Optional last paragraph: the tools used (PyTorch \citep{paszke2019pytorch},
MLflow, Optuna \citep{akiba2019optuna}, Git, \LaTeX{}) and the environment differences
between my local Windows machine and the Grid'5000 Linux nodes.}

{The stack is PyTorch \citep{paszke2019pytorch} for the models, MLflow for experiment
tracking, Optuna \citep{akiba2019optuna} for the hyperparameter searches referenced in
\Cref{app:repro}, Git for version control, and \LaTeX{} for this report. Day-to-day
development and small smoke tests happened locally on Windows; every real training run
happened on Grid'5000 Linux GPU nodes (\Cref{tab:protocol}). The two environments
differ in more than the operating system -- Python and package versions, and how a job
is started (a local terminal versus an \texttt{oarsub} submission) -- so the same
\texttt{configs/*.yaml} files are used on both sides to keep the one thing that must
not drift (the model and training recipe) identical between them.}